%% Dans ce fichier, les �ventuels accents sont cod�s
%% en ISO-8859-15 (aka latin9)
%% La ligne suivante donne trois exemples 
%% e accent aigu : �, e accent grave : �, e-dans-l'o : �
%%%%%ENTREES%%%%
\begin{tabarg}
\multicolumn{3}{l}{\titretab {Entr\'ees}}\\
\hline
$X\_in$ & Vecteur, \texttt{float} &  Vecteur d'entr�e\\
$P\_in$ & Vecteur, \texttt{float} & Altitude en pression\\
$tgt\_lat$ & Coeff., \texttt{float} & Latitude du point cible en degr� d�cimal\\
$tgt\_lon$ & Coeff., \texttt{float} & Longitude du point cible en degr� d�cimal\\
$tgt\_alt$ & Coeff., \texttt{float} & Altitude du point cible en degr� d�cimal\\
\end{tabarg}

\medskip
%%%%%SORTIES%%%%
\begin{tabarg}
\multicolumn{3}{l}{\titretab {Sorties}}\\
\hline
$X\_out$ & Vecteur, \texttt{float} & Colonne extraite du vecteur d'entr�e la plus proche du point cible\\
\end{tabarg}
